\documentclass[11pt,a4paper]{article}

\usepackage[utf8]{inputenc}
\usepackage[T1]{fontenc}
\usepackage{lmodern}
\usepackage{amsmath,amssymb,amsfonts,bm}
\usepackage{geometry}
\usepackage{hyperref}
\usepackage{authblk}
\usepackage{booktabs}
\usepackage{enumitem}

\geometry{margin=2.4cm}

\title{\textbf{Technical Note (v49.1)}\\[0.2cm]
\textbf{Effective Gravitational Coupling in RDCC}\\[0.15cm]
\large Reconciliation of \(G_{\mathrm{eff}}\) between Companions I and XI}

\author[1]{Michael Lehmann}
\affil[1]{\small Independent Researcher, Relaxation-Driven Cyclic Cosmology (RDCC)\\
\texttt{mi.lehmann@gmx.de}}

\date{September 2026}

\begin{document}
\maketitle

\begin{abstract}
Companions I and XI of RDCC v49.1 employ two incompatible expressions for the
effective Newton constant inferred by a visible-sector observer. Companion I
obtains
\(G_{\mathrm{eff}}=G_N\bigl(1+\rho_{(-)}/\rho_{(+)}\bigr)\)
from equating the global and sectoral linearized Einstein responses.
Companion XI writes the reciprocal, including a relaxation piece in the
denominator. Both papers share the same global Friedmann equation
\(3H^2=8\pi G_N(\rho_{(+)}+\rho_{(-)}+\rho_{\mathrm{rel}})\). That equation
fixes the projection uniquely: the Companion-I form is the consistent
result, with an optional \(\mathcal{O}(\alpha_{\mathrm{IR}})\) extension
\(\rho_{\mathrm{rel}}/\rho_{(+)}\). The working \(\Omega\)-form of \(H(z)\)
in Companion XI remains valid and does not require the inverted
\(G_{\mathrm{eff}}\). This note records the algebra and the replacement rule
for Companion XI. It introduces no new parameter and no new interaction.
\end{abstract}

\tableofcontents
\newpage

\section{Purpose}

This note consolidates one structural inconsistency in RDCC v49.1. It does
not modify the bipartite CPT postulate, the dimension-8 operator, or the
freeze-out scale of Companion XII. It answers a single question: which
expression for \(G_{\mathrm{eff}}\) follows from the global Einstein
equation already adopted in both Companions I and XI?

\section{The two published formulae}

Companion I derives the sectoral coupling from
\begin{equation}
  \delta G_{\mu\nu}
  =
  8\pi G_N
  \bigl(
    \delta T^{(+)}_{\mu\nu}
    +
    \delta T^{(-)}_{\mu\nu}
  \bigr)
\end{equation}
and from the sectoral rewriting
\begin{equation}
  \delta G_{\mu\nu}
  =
  8\pi G_{\mathrm{eff}}\,
  \delta T^{(+)}_{\mu\nu}.
\end{equation}
Equating the two responses yields
\begin{equation}
  G_{\mathrm{eff}}
  =
  G_N
  \left(
    1
    +
    \frac{\delta T^{(-)}_{\mu\nu}}{\delta T^{(+)}_{\mu\nu}}
  \right).
  \label{eq:Ilinear}
\end{equation}
On a homogeneous FRW background this becomes
\begin{equation}
  G_{\mathrm{eff}}^{\mathrm{(I)}}
  =
  G_N
  \left(
    1
    +
    \frac{\rho_{(-)}}{\rho_{(+)}}
  \right).
  \label{eq:GI}
\end{equation}

Companion XI instead writes
\begin{equation}
  G_{\mathrm{eff}}^{\mathrm{(XI)}}
  =
  \frac{G_N}{
    1
    +
    \rho_{(-)}/\rho_{(+)}
    +
    \rho_{\mathrm{rel}}/\rho_{(+)}
  }.
  \label{eq:GXI}
\end{equation}
Equations~\eqref{eq:GI} and~\eqref{eq:GXI} are not equivalent. At most one
of them can be the projection of the same global Einstein equation.

\section{Global constraint}

Both Companions accept
\begin{equation}
  G_{\mu\nu}
  =
  8\pi G_N
  \bigl(
    T^{(+)}_{\mu\nu}
    +
    T^{(-)}_{\mu\nu}
    +
    T^{\mathrm{rel}}_{\mu\nu}
  \bigr)
\end{equation}
and the background reduction
\begin{equation}
  3H^2
  =
  8\pi G_N
  \bigl(
    \rho_{(+)}
    +
    \rho_{(-)}
    +
    \rho_{\mathrm{rel}}
  \bigr).
  \label{eq:globalH}
\end{equation}
A visible-sector observer who parametrizes the same expansion as
\begin{equation}
  3H^2
  =
  8\pi G_{\mathrm{eff}}\, \rho_{(+)}
  +
  \Lambda_{\mathrm{eff}}
  \label{eq:sectorH}
\end{equation}
must match coefficients with Eq.~\eqref{eq:globalH}. Absorbing
\(\rho_{\mathrm{rel}}\) either into \(G_{\mathrm{eff}}\) or into
\(\Lambda_{\mathrm{eff}}\) is a bookkeeping choice; it cannot invert the
sum \(\rho_{(+)}+\rho_{(-)}\).

Identifying the matter piece of Eq.~\eqref{eq:sectorH} with the first two
terms of Eq.~\eqref{eq:globalH} gives immediately
\begin{equation}
  G_{\mathrm{eff}}
  =
  G_N
  \left(
    1
    +
    \frac{\rho_{(-)}}{\rho_{(+)}}
    +
    \frac{\rho_{\mathrm{rel}}}{\rho_{(+)}}
  \right),
  \label{eq:Gfinal}
\end{equation}
or, if \(\rho_{\mathrm{rel}}\) is kept inside \(\Lambda_{\mathrm{eff}}\),
the Companion-I result Eq.~\eqref{eq:GI}. In either convention
\(G_{\mathrm{eff}}\ge G_N\). The reciprocal Eq.~\eqref{eq:GXI} would make
gravity \emph{weaker} as the mirror density grows, which contradicts
Eq.~\eqref{eq:globalH}: adding \(\rho_{(-)}\) increases \(H\), it does not
decrease \(G_N\rho_{(+)}\).

\section{Adopted definition}

The Phase-1 definition is therefore
\begin{equation}
  \boxed{
    G_{\mathrm{eff}}
    =
    G_N
    \left(
      1
      +
      \frac{\rho_{(-)}}{\rho_{(+)}}
      +
      \frac{\rho_{\mathrm{rel}}}{\rho_{(+)}}
    \right).
  }
  \label{eq:box}
\end{equation}
On backgrounds where \(\rho_{\mathrm{rel}}/\rho_{(+)}\) is of order
\(\alpha_{\mathrm{IR}}\simeq0.018\), the last term may be dropped and
Eq.~\eqref{eq:GI} recovered. Companion I is the correct leading projection.
Companion XI must replace Eq.~\eqref{eq:GXI} by Eq.~\eqref{eq:box}.

\section{What does not change}

The practical Hubble function of Companion XI,
\begin{equation}
  H^2(z)
  =
  H_0^2
  \Bigl[
    \Omega_b(1+z)^3
    +
    \Omega_{\mathrm{DM,eff}}(1+z)^3
    +
    \Omega_r(1+z)^4
    +
    \Omega_\Lambda^{\mathrm{eff}}
    +
    \delta_{\mathrm{rel}}(z)
  \Bigr],
\end{equation}
already implements Eq.~\eqref{eq:globalH} in \(\Omega\)-language:
\(\Omega_{\mathrm{DM,eff}}\) counts \(\rho_{(-)}\) with the same \(G_N\)
that multiplies \(\rho_{(+)}\). That equation is consistent and should be
kept. It must not be rewritten with the inverted \(G_{\mathrm{eff}}\),
which would double-count or cancel the mirror contribution.

Linear growth and gravitational slip continue to use the global Poisson
structure of Companion VII,
\begin{equation}
  \nabla^2\Phi-3H(\dot\Phi+H\Phi)
  =
  4\pi G_N
  \bigl(
    \rho_{(+)}\delta_{(+)}
    +
    \rho_{(-)}\delta_{(-)}
  \bigr)
  +
  \cdots,
\end{equation}
not a sectoral \(G_{\mathrm{eff}}\) inserted only on \(\delta_{(+)}\).
The background projection Eq.~\eqref{eq:box} and the perturbation source
are the same global coupling, written once for FRW and once for
\(\delta\rho\).

\section{Replacement rule for Companion XI}

In the subsection ``Visible-sector projection'' of Companion XI, replace
Eq.~\eqref{eq:GXI} by Eq.~\eqref{eq:box} and add:
\begin{quote}
The inverted expression published in an earlier draft of this Companion
is incompatible with the global Friedmann constraint
\(3H^2=8\pi G_N(\rho_{(+)}+\rho_{(-)}+\rho_{\mathrm{rel}})\). The
\(\Omega\)-form of \(H(z)\) is unchanged.
\end{quote}
No other companion that writes \(G_{\mathrm{eff}}=G_N(1+\rho_-/\rho_+)\)
requires amendment. Any occurrence of the reciprocal should be treated as
superseded by this note.

\section{Relation to the \texorpdfstring{\(k_{\mathrm{rel}}\)}{k\_rel} note}

The numerical evaluation of
\(k_{\mathrm{rel}}=a_{\mathrm{fr}}H(a_{\mathrm{fr}})\) uses the
\(\Omega\)-form of \(H(a)\), hence \(G_N\) and both densities, not the
inverted \(G_{\mathrm{eff}}\). That table is therefore stable under the
present reconciliation. What would have been unstable is a background
integrated with Eq.~\eqref{eq:GXI}: mirror density would have suppressed
rather than sourced expansion.

\section{What this note does not decide}

The note does not derive \(\rho_{(-)}/\rho_{(+)}\) from first principles,
does not identify that ratio with \(\Omega_{\mathrm{DM}}/\Omega_b\), and
does not promote \(\rho_{\mathrm{rel}}\) beyond an
\(\mathcal{O}(\alpha_{\mathrm{IR}})\) bookkeeping term. Those are separate
questions. The only claim is the projection identity
Eq.~\eqref{eq:box}.

\end{document}